
\documentclass{article}
\usepackage{colortbl}
\usepackage{makecell}
\usepackage{multirow}
\usepackage{supertabular}

\begin{document}

\newcounter{utterance}

\centering \large Interaction Transcript for game `cladder', experiment `full\_v1.5\_default', episode 2803 with reflexion{-}Qwen3.5{-}27B.
\vspace{24pt}

{ \footnotesize  \setcounter{utterance}{1}
\setlength{\tabcolsep}{0pt}
\begin{supertabular}{c@{$\;$}|p{.15\linewidth}@{}p{.15\linewidth}p{.15\linewidth}p{.15\linewidth}p{.15\linewidth}p{.15\linewidth}}
    \# & \multicolumn{2}{c}{Player} && \multicolumn{2}{c}{Game Master} \\
    \hline

    \theutterance \stepcounter{utterance}  
    & & & \multicolumn{4}{p{0.6\linewidth}}{
        \cellcolor[rgb]{0.9,0.9,0.9}{
            \makecell[{{p{\linewidth}}}]{
                \texttt{\tiny{[P1$\langle$GM]}}
                \texttt{You have attempted this same task before and did not succeed. These are the plans you wrote after those attempts.} \\
\texttt{Trial \#1: New plan: My previous attempt failed because I over{-}analyzed the logical gap between "not dark" and "bright" and, more critically, I failed to provide a complete, definitive answer, cutting off mid{-}sentence. In this specific causal inference context, the prompt establishes a binary world where the only defined mechanism for darkness is the conjunction of "blowing out" and "having wax." Since the candle lacks wax, the intervention of blowing it out cannot trigger the "dark" state. Therefore, the room remains in its default state (bright). Next time, I will immediately conclude that the room is bright because the necessary condition for darkness (wax) is absent, rendering the action of blowing out the candle ineffective. I will ensure the response starts with a clear "Yes" or "No" and provides a concise, complete explanation without trailing off or introducing undefined states.} \\
\\ 
\\ 
\texttt{You are an expert in causal inference. The following question is not a typical commonsense query, but rather a meticulously designed question created by a professor specializing in causal inference, intended to assess the students' mastery of the course content.} \\
\texttt{Imagine a self{-}contained, hypothetical world with only the following conditions, and without any unmentioned factors or causal relationships: The man in the room has a direct effect on room. The candle has a direct effect on room. We know that blowing out the candle and candle with wax causes dark room. We observed the candle is out of wax. Would the room is bright if blowing out the candle instead of not blowing out the candle?} \\
\texttt{Start your answer with 'Yes' or 'No', followed by additional reasoning or evidence to support your explanation.} \\
            }
        }
    }
     \\ \\

    \theutterance \stepcounter{utterance}  
    & \multicolumn{4}{p{0.6\linewidth}}{
        \cellcolor[rgb]{0.561,0.737,0.561}{
            \makecell[{{p{\linewidth}}}]{
                \texttt{\tiny{[P1$\rangle$GM]}}
                \texttt{Yes} \\
\\ 
\texttt{In this strictly defined hypothetical world, the causal mechanism for a "dark room" is explicitly stated as the conjunction of two conditions: "blowing out the candle" AND "candle with wax." This establishes a logical structure where Darkness = (Blow Out) ∧ (Has Wax).} \\
\\ 
\texttt{The scenario specifies that the candle is "out of wax," meaning the condition (Has Wax) is false. In causal logic, if a necessary condition for an effect is absent, the effect cannot occur, regardless of the status of other variables. Therefore, even if the action "blowing out the candle" is performed, the conjunction required to produce darkness fails because the wax condition is not met.} \\
\\ 
\texttt{Since the only defined mechanism for darkness is broken, the room cannot become dark. In the absence of any other factors or causal relationships that could alter the room's state, the room remains in its default state, which is bright. Thus, blowing out the candle (in the absence of wax) results in a bright room, just as not blowing it out would.} \\
            }
        }
    }
    & & \\ \\

    \theutterance \stepcounter{utterance}  
    & & & \multicolumn{2}{p{0.3\linewidth}}{
        \cellcolor[rgb]{0.9,0.9,0.9}{
            \makecell[{{p{\linewidth}}}]{
                \texttt{\tiny{[GM$|$GM]}}
                \texttt{Yes} \\
            }
        }
    }
    & & \\ \\

    \theutterance \stepcounter{utterance}  
    & & & \multicolumn{2}{p{0.3\linewidth}}{
        \cellcolor[rgb]{0.9,0.9,0.9}{
            \makecell[{{p{\linewidth}}}]{
                \texttt{\tiny{[GM$|$GM]}}
                \texttt{yes} \\
            }
        }
    }
    & & \\ \\

    \theutterance \stepcounter{utterance}  
    & & & \multicolumn{2}{p{0.3\linewidth}}{
        \cellcolor[rgb]{0.9,0.9,0.9}{
            \makecell[{{p{\linewidth}}}]{
                \texttt{\tiny{[GM$|$GM]}}
                \texttt{game\_result = WIN} \\
            }
        }
    }
    & & \\ \\

\end{supertabular}
}

\end{document}
